% ============================================================
%  Rapport de stage — Sections Encodage / Log / Standardisation
%  À inclure dans votre mémoire (chapitre 3).
%  Remplacez X par le numéro de section approprié.
% ============================================================

\documentclass[a4paper,12pt]{article}
\usepackage[utf8]{inputenc}
\usepackage[T1]{fontenc}
\usepackage[french]{babel}
\usepackage{amsmath, amssymb}
\usepackage{booktabs}
\usepackage{geometry}
\geometry{margin=2.5cm}

\begin{document}

% ------------------------------------------------------------
\section{Encodage des variables}
% ------------------------------------------------------------

Une fois les variables brutes nettoyées et les nouvelles caractéristiques construites, il est indispensable de transformer les variables catégorielles et temporelles en représentations numériques exploitables par les algorithmes de machine learning. Le pipeline applique ainsi un encodage cyclique pour la variable temporelle, un encodage par rareté pour les variables catégorielles, complétés par une transformation logarithmique et une standardisation.

\subsection{Encodage cyclique de l'heure}

L'heure est une variable temporelle de nature périodique : 23~h et 0~h sont deux valeurs très proches sur le cycle d'une journée, alors qu'une simple numérotation les éloignerait artificiellement (23 et 0). Pour préserver cette proximité, nous avons appliqué un encodage cyclique basé sur les fonctions trigonométriques :
\begin{equation}
  \texttt{date\_hour\_sin} = \sin\!\left(\frac{2\pi \times \text{heure}}{24}\right), \qquad
  \texttt{date\_hour\_cos} = \cos\!\left(\frac{2\pi \times \text{heure}}{24}\right).
\end{equation}
Cet encodage projette l'heure sur un cercle unité : chaque heure est ainsi représentée par une coordonnée bidimensionnelle unique (sin, cos), et la distance entre deux heures tient naturellement compte de la périodicité. Cette approche évite les discontinuités induites par un encodage ordinal classique. De plus, le temps écoulé depuis la première requête (\texttt{date\_normalized}) a été calculé en secondes afin de quantifier l'ordre chronologique des requêtes dans l'analyse.

\subsection{Encodage par rareté (Rarity Encoding)}

Les variables catégorielles (méthode HTTP, catégorie de statut) ont été encodées selon leur fréquence d'apparition. La fonction d'encodage par rareté associe à chaque catégorie un score défini comme :
\begin{equation}
  \text{score} = 1 - \text{fréquence relative}(\text{catégorie}), \qquad
  \text{fréquence relative} = \frac{\text{nombre d'occurrences de la catégorie}}{\text{nombre total d'observations}}.
\end{equation}
Ainsi, une catégorie fréquente obtient un score proche de 0, tandis qu'une catégorie rare obtient un score proche de 1.

Cette méthode présente plusieurs avantages : elle n'augmente pas la dimensionnalité du jeu de données (contrairement au one-hot encoding qui aurait créé autant de colonnes que de catégories) ; elle conserve l'information d'apparition des catégories tout en restant linéaire en temps de calcul ; et elle est robuste face aux catégories nombreuses ou à forte cardinalité.

En pratique, sur la variable \texttt{method}, le score obtenu varie entre \textbf{0{,}1095} (méthode GET, très fréquente : fréquence relative de 0{,}891) et \textbf{0{,}9961} (méthode POST, rare : 0{,}004), pour une moyenne de \textbf{0{,}1958}. Concernant la catégorie de statut HTTP, le score est nul pour toutes les observations : dans notre jeu de données, toutes les requêtes appartiennent à la même catégorie (codes 301/302 $\to$ classe \og 3xx \fg{}), sa fréquence relative vaut donc 1 et le score $1 - 1 = 0$. Cette colonne étant constante, elle s'avérera non informative pour la modélisation et pourra être écartée lors de la sélection de variables.

% ------------------------------------------------------------
\section{Transformation logarithmique}
% ------------------------------------------------------------

Afin de réduire l'asymétrie des distributions fortement étalées (taille de réponse, nombre de requêtes par IP, profondeur d'URL, etc.), une transformation logarithmique de type $\log(1 + x)$ a été appliquée sur \textbf{19 variables}. L'ajout de la constante 1 permet de conserver les valeurs nulles dans le domaine de définition.

\subsection{Définition d'une distribution fortement asymétrique}

Une distribution est dite fortement asymétrique (ou \emph{skewed}) lorsque ses valeurs ne sont pas réparties symétriquement autour de la moyenne, mais présentent une queue longue d'un côté. Dans le cas d'une asymétrie positive (queue droite), la majorité des valeurs sont petites, tandis que quelques valeurs extrêmes sont très grandes. Cela se traduit par une médiane inférieure à la moyenne et un écart-type élevé par rapport à la plage typique des données. Par exemple, une variable comme \texttt{ip\_request\_count}, avec une médiane de 36 et un maximum de 1~411, illustre cette propriété. À l'inverse, une distribution symétrique (ex. normale) a des queues équilibrées, comme observé pour \texttt{date\_hour\_sin}.

\subsection{Impact sur la détection des intrusions}

Dans le contexte de la détection des intrusions, une distribution asymétrique peut poser plusieurs problèmes :
\begin{itemize}
  \item \textbf{Biais dans les modèles :} Les modèles non supervisés comme PCA ou Isolation Forest s'appuient sur des distances ou des variances. Les valeurs extrêmes dans une queue droite (ex. \texttt{ip\_interval\_std} avec un maximum de 1~885~947 secondes) dominent ces calculs, masquant les variations subtiles des petites valeurs qui pourraient indiquer des anomalies (ex. bursts ou injections).
  \item \textbf{Difficulté à identifier les anomalies :} Les comportements malveillants (ex. scans avec des centaines de requêtes par IP) se situent souvent dans les queues longues, mais leur impact est exagéré sans transformation, rendant les seuils de détection moins discriminants.
  \item \textbf{Influence disproportionnée des extrêmes :} Une variable non transformée comme \texttt{ip\_request\_rate} (max $= 79$, médiane $= 0{,}0003$) peut biaiser les composantes principales ou les scores d'anomalie, au détriment des variables plus symétriques.
\end{itemize}
Ces effets sont amplifiés lorsque les variables ont des échelles très différentes, un défi déjà abordé par la standardisation (Section \ref{sec:standardisation}), mais insuffisant pour corriger l'asymétrie.

\section{Solution : transformation logarithmique}

La transformation $\log(1 + x)$ est appliquée aux variables asymétriques avant standardisation. Cette méthode :
\begin{itemize}
  \item \textbf{Réduit l'asymétrie :} Elle compresse les grandes valeurs et amplifie les différences entre petites valeurs. Par exemple, $\log(1 + 1\,411) \approx 7{,}25$ contre $\log(1 + 36) \approx 3{,}61$ pour \texttt{ip\_request\_count}, rendant la distribution plus symétrique (le coefficient d'asymétrie passe de 1,89 à 0,02).
  \item \textbf{Améliore la détection :} Les anomalies dans les petites valeurs (ex. taux de requêtes faibles mais suspects) deviennent plus visibles, tandis que les extrêmes sont moins dominants.
  \item \textbf{Prépare les modèles :} Une distribution plus proche de la normale facilite l'interprétation des variances dans PCA ou des distances dans Isolation Forest.
\end{itemize}
Cependant, cette transformation n'est pas systématique : pour les variables symétriques ou binaires, elle est inutile voire contre-productive, car elle altère leur granularité naturelle.

% ------------------------------------------------------------
\section{Analyse des variables pour la transformation logarithmique}
% ------------------------------------------------------------

Cette section examine les variables du jeu de données par catégorie, en se basant sur leurs statistiques descriptives (moyenne, écart-type, min, 25~\%~, 50~\%, 75~\%, max) calculées sur les \textbf{41~433 lignes} du jeu de données \texttt{Access\_log\_cleaned.csv}, pour identifier celles qui bénéficient d'une transformation logarithmique.

\subsection{Variables binaires}

\begin{itemize}
  \item \textbf{\texttt{is\_bot} :} Moyenne $= 0{,}3777$, Écart-type $= 0{,}4848$, Min $= 0$, 25~\% $= 0$, 50~\% $= 0$, 75~\% $= 1$, Max $= 1$. \emph{Observation :} Distribution binaire (0 ou 1), où 37{,}8~\% des requêtes proviennent de bots. Aucune transformation logarithmique n'est applicable ni nécessaire.
  \item \textbf{\texttt{referer\_is\_present} :} Moyenne $= 0{,}0863$, Écart-type $= 0{,}2808$, Min $= 0$, 50~\% $= 0$, 75~\% $= 0$, Max $= 1$. \emph{Observation :} Distribution binaire avec faible proportion de 1 (8{,}6~\%). Pas de transformation logarithmique nécessaire.
  \item \textbf{\texttt{is\_hierarchical} :} Moyenne $= 0{,}4128$, Écart-type $= 0{,}4923$, Min $= 0$, 50~\% $= 0$, 75~\% $= 1$, Max $= 1$. \emph{Observation :} Distribution binaire. Pas de log applicable ni utile.
\end{itemize}

\subsection{Variables temporelles}

\begin{itemize}
  \item \textbf{\texttt{date\_hour\_sin} :} Moyenne $= -0{,}0875$, Écart-type $= 0{,}7031$, Min $= -1$, 25~\% $= -0{,}7071$, 50~\% $= -0{,}2588$, 75~\% $= 0{,}5000$, Max $= 1$. \emph{Observation :} Distribution cyclique symétrique entre $-1$ et $1$. Pas d'asymétrie ; log inutile.
  \item \textbf{\texttt{date\_hour\_cos} :} Moyenne $= -0{,}0727$, Écart-type $= 0{,}7020$, Min $= -1$, 25~\% $= -0{,}7071$, 50~\% $= -0{,}2588$, 75~\% $= 0{,}7071$, Max $= 1$. \emph{Observation :} Distribution cyclique symétrique. Pas de log nécessaire.
  \item \textbf{\texttt{date\_normalized} :} Moyenne $= 1~370~392$, Écart-type $= 783~027$, Min $= 0$, 25~\% $= 702~360$, 50~\% $= 1~379~361$, 75~\% $= 2~058~480$, Max $= 2~782~661$. \emph{Observation :} Distribution linéaire homogène sur la période d'observation ($\sim$32 jours). La forme brute préserve la granularité pour détecter les bursts ; pas de log nécessaire.
\end{itemize}

\subsection{Variables liées à l'IP}

\begin{itemize}
  \item \textbf{\texttt{ip\_request\_count} :} Moyenne $= 233{,}3$, Écart-type $= 353{,}5$, Min $= 1$, 25~\% $= 11$, 50~\% $= 36$, 75~\% $= 403$, Max $= 1~411$. \emph{Observation :} Forte asymétrie (médiane $\ll$ moyenne, max élevé). Queue droite ; \textbf{log utile}.
  \item \textbf{\texttt{ip\_time\_spread} :} Moyenne $= 1~008~131$, Écart-type $= 1~014~842$, Min $= 0$, 25~\% $= 80~719$, 50~\% $= 516~793$, 75~\% $= 2~037~141$, Max $= 2~780~186$. \emph{Observation :} Asymétrie marquée (médiane $<$ moyenne). Queue droite ; \textbf{log utile}.
  \item \textbf{\texttt{ip\_subnet\_activity} :} Moyenne $= 16{,}5$, Écart-type $= 16{,}5$, Min $= 1$, 25~\% $= 2$, 50~\% $= 10$, 75~\% $= 26$, Max $= 75$. \emph{Observation :} Distribution relativement symétrique (médiane proche de la moyenne). Pas de log nécessaire.
  \item \textbf{\texttt{ip\_error\_rate} :} Moyenne $= 0$, Écart-type $= 0$, Min $= 0$, Max $= 0$. \emph{Observation :} Constante à 0 (aucun code $\geq 400$ dans les données). Variable non informative, à écarter.
  \item \textbf{\texttt{ip\_request\_rate} :} Moyenne $= 0{,}3394$, Écart-type $= 3{,}5081$, Min $= 7{,}5\times 10^{-7}$, 25~\% $= 0{,}0001$, 50~\% $= 0{,}0003$, 75~\% $= 0{,}0005$, Max $= 79$. \emph{Observation :} Extrêmement asymétrique (médiane très basse, max élevé), plusieurs ordres de grandeur (coefficient d'asymétrie de 21,5 avant log, 5,4 après). \textbf{Log fortement recommandé}.
  \item \textbf{\texttt{ip\_interval\_std} :} Moyenne $= 35~919$, Écart-type $= 97~903$, Min $= 0$, 25~\% $= 2~105$, 50~\% $= 5~702$, 75~\% $= 21~755$, Max $= 1~885~947$. \emph{Observation :} Forte asymétrie (médiane $<$ moyenne, max extrême). Queue droite ; \textbf{log utile}.
  \item \textbf{\texttt{ip\_method\_diversity} :} Moyenne $= 1{,}0666$, Écart-type $= 0{,}2493$, Min $= 1$, 50~\% $= 1$, 75~\% $= 1$, Max $= 2$. \emph{Observation :} Plage très limitée (1--2). Pas de log nécessaire.
\end{itemize}

\subsection{Variables liées à l'URL}

\begin{itemize}
  \item \textbf{\texttt{nbre\_unique\_urls\_per\_ip} :} Moyenne $= 36{,}8$, Écart-type $= 48{,}8$, Min $= 1$, 25~\% $= 1$, 50~\% $= 13$, 75~\% $= 65$, Max $= 198$. \emph{Observation :} Forte asymétrie (médiane $\ll$ moyenne). Queue droite ; \textbf{log utile}.
  \item \textbf{\texttt{url\_path\_struct\_entropy} :} Moyenne $= 5{,}17$, Écart-type $= 6{,}38$, Min $= 0$, 25~\% $= 0$, 50~\% $= 2{,}65$, 75~\% $= 12{,}99$, Max $= 22{,}75$. \emph{Observation :} Asymétrie (max élevé). Queue droite ; \textbf{log utile}.
  \item \textbf{\texttt{url\_path\_depth} :} Moyenne $= 2{,}07$, Écart-type $= 2{,}58$, Min $= 0$, 25~\% $= 0$, 50~\% $= 1$, 75~\% $= 5$, Max $= 8$. \emph{Observation :} Asymétrie (75~\% $= 5$, max élevé). Queue droite ; \textbf{log utile}.
  \item \textbf{\texttt{nbre\_total\_url\_query\_params} :} Moyenne $= 0{,}26$, Écart-type $= 0{,}69$, Min $= 0$, 50~\% $= 0$, 75~\% $= 0$, Max $= 7$. \emph{Observation :} Très asymétrique (75~\% $= 0$). Queue droite ; \textbf{log utile}.
  \item \textbf{\texttt{url\_query\_special\_chars\_count} :} Moyenne $= 2{,}46$, Écart-type $= 6{,}03$, Min $= 0$, 50~\% $= 0$, 75~\% $= 0$, Max $= 89$. \emph{Observation :} Extrêmement asymétrique (75~\% $= 0$, max très élevé). Queue droite ; \textbf{log utile}.
  \item \textbf{\texttt{url\_path\_special\_chars\_count} :} Moyenne $= 3{,}29$, Écart-type $= 2{,}91$, Min $= 1$, 25~\% $= 1$, 50~\% $= 2$, 75~\% $= 7$, Max $= 27$. \emph{Observation :} Asymétrie (75~\% $= 7$, max élevé). Queue droite ; \textbf{log utile}.
  \item \textbf{\texttt{fragment\_param\_count} et \texttt{special\_char\_fragment\_count} :} Moyenne $= 0$, Écart-type $= 0$, Min $= 0$, Max $= 0$. \emph{Observation :} Constantes à 0 dans ce jeu de données (aucune URL avec fragment). Non informatives.
\end{itemize}

\subsection{Variables liées au User-Agent}

\begin{itemize}
  \item \textbf{\texttt{user\_agent\_length} :} Moyenne $= 87{,}3$, Écart-type $= 43{,}8$, Min $= 0$, 25~\% $= 71$, 50~\% $= 72$, 75~\% $= 115$, Max $= 306$. \emph{Observation :} Légère asymétrie (max élevé), mais distribution relativement symétrique. Pas de log nécessaire.
  \item \textbf{\texttt{user\_agent\_entropy} :} Moyenne $= 4{,}40$, Écart-type $= 1{,}37$, Min $= 0$, 25~\% $= 4{,}49$, 50~\% $= 4{,}71$, 75~\% $= 5{,}18$, Max $= 5{,}48$. \emph{Observation :} Symétrique (plage limitée). Pas d'asymétrie ; pas de log nécessaire.
  \item \textbf{\texttt{user\_agent\_total\_special\_chars} :} Moyenne $= 15{,}0$, Écart-type $= 6{,}3$, Min $= 0$, 25~\% $= 12$, 50~\% $= 16$, 75~\% $= 18$, Max $= 50$. \emph{Observation :} Légère asymétrie, mais pas critique. Pas de log nécessaire.
\end{itemize}

\subsection{Variables liées au Referer}

\begin{itemize}
  \item \textbf{\texttt{referer\_path\_struct\_entropy} :} Moyenne $= 0{,}028$, Écart-type $= 0{,}38$, Min $= 0$, 50~\% $= 0$, 75~\% $= 0$, Max $= 12{,}92$. \emph{Observation :} Très asymétrique (75~\% $= 0$). Queue droite ; \textbf{log utile}.
  \item \textbf{\texttt{referer\_path\_depth} :} Moyenne $= 0{,}0097$, Écart-type $= 0{,}139$, Min $= 0$, 50~\% $= 0$, 75~\% $= 0$, Max $= 5$. \emph{Observation :} Très asymétrique (75~\% $= 0$). Queue droite ; \textbf{log utile}.
  \item \textbf{\texttt{nbre\_total\_referer\_query\_params} :} Moyenne $= 0{,}0055$, Écart-type $= 0{,}175$, Min $= 0$, 50~\% $= 0$, 75~\% $= 0$, Max $= 10$. \emph{Observation :} Très asymétrique (75~\% $= 0$). Queue droite ; \textbf{log utile}.
  \item \textbf{\texttt{referer\_query\_special\_chars\_count} :} Moyenne $= 0{,}021$, Écart-type $= 0{,}711$, Min $= 0$, 50~\% $= 0$, 75~\% $= 0$, Max $= 59$. \emph{Observation :} Très asymétrique (75~\% $= 0$). Queue droite ; \textbf{log utile}.
\end{itemize}

\subsection{Variable liée à la taille de réponse}

\begin{itemize}
  \item \textbf{\texttt{response\_size} :} Moyenne $= 232{,}8$, Écart-type $= 94{,}1$, Min $= 0$, 25~\% $= 214$, 50~\% $= 234$, 75~\% $= 270$, Max $= 892$. \emph{Observation :} Distribution concentrée (75~\% $\leq 270$) mais présentant des valeurs nulles et un maximum de 892 octets ; \textbf{log utile} pour réduire l'impact des valeurs extrêmes.
\end{itemize}

\subsection{Variables encodées}

\begin{itemize}
  \item \textbf{\texttt{method\_encoded} :} Moyenne $= 0{,}1958$, Écart-type $= 0{,}2463$, Min $= 0{,}1095$, 50~\% $= 0{,}1095$, 75~\% $= 0{,}1095$, Max $= 0{,}9961$. \emph{Observation :} Asymétrie (75~\% $= 0{,}1095$, max proche de 1). Queue droite ; \textbf{log utile}.
  \item \textbf{\texttt{status\_code\_category\_encoded} :} Moyenne $= 0$, Écart-type $= 0$, Min $= 0$, Max $= 0$. \emph{Observation :} Constante à 0 (encodage dégénéré, cf. Section sur l'encodage). Pas de log applicable ; variable à écarter.
\end{itemize}

\subsection{Synthèse}

Sur les variables analysées, \textbf{19} présentent une asymétrie suffisante pour bénéficier d'une transformation logarithmique $\log(1 + x)$ : \texttt{response\_size}, \texttt{ip\_request\_count}, \texttt{ip\_time\_spread}, \texttt{ip\_request\_rate}, \texttt{ip\_interval\_std}, \texttt{nbre\_unique\_urls\_per\_ip}, \texttt{url\_path\_struct\_entropy}, \texttt{url\_path\_depth}, \texttt{nbre\_total\_url\_query\_params}, \texttt{url\_query\_special\_chars\_count}, \texttt{url\_path\_special\_chars\_count}, \texttt{referer\_path\_struct\_entropy}, \texttt{referer\_path\_depth}, \texttt{referer\_query\_special\_chars\_count}, \texttt{nbre\_total\_referer\_query\_params}, \texttt{method\_encoded}, ainsi que \texttt{ip\_error\_rate}, \texttt{fragment\_param\_count} et \texttt{special\_char\_fragment\_count} (incluses dans le pipeline mais constantes à 0 dans ce jeu de données, donc non informatives). Les variables restantes, symétriques, binaires ou à plage limitée (\texttt{is\_bot}, \texttt{referer\_is\_present}, \texttt{is\_hierarchical}, \texttt{date\_hour\_sin}, \texttt{date\_hour\_cos}, \texttt{date\_normalized}, \texttt{ip\_subnet\_activity}, \texttt{ip\_method\_diversity}, \texttt{user\_agent\_length}, \texttt{user\_agent\_entropy}, \texttt{user\_agent\_total\_special\_chars}, \texttt{status\_code\_category\_encoded}), n'en ont pas besoin. Cette approche, combinée à la standardisation (Section \ref{sec:standardisation}), optimise la détection en uniformisant les distributions pour les modèles non supervisés.

% ------------------------------------------------------------
\section{Standardisation}
\label{sec:standardisation}
% ------------------------------------------------------------

Après l'encodage et la transformation logarithmique, une standardisation a été appliquée à l'ensemble des variables numériques continues afin de les rendre comparables. Cette étape, réalisée à l'aide du \texttt{StandardScaler} de scikit-learn, transforme chaque variable selon la formule :
\begin{equation}
  x_{\text{scaled}} = \frac{x - \text{moyenne}}{\text{écart-type}}
\end{equation}
où la moyenne et l'écart-type sont ceux de la variable. Le résultat est une distribution de \textbf{moyenne nulle} ($\mu = 0$) et \textbf{d'écart-type unitaire} ($\sigma = 1$) pour chaque variable standardisée.

Ont été exclues de cette transformation : les variables binaires (\texttt{is\_bot}, \texttt{is\_hierarchical}, \texttt{referer\_is\_present}), les colonnes d'identification et de type (\texttt{status\_code}, \texttt{pid}, \texttt{tid}) et la colonne temporelle \texttt{datetime}. Au total, \textbf{57 colonnes} ont été standardisées.

Cette étape est indispensable pour plusieurs raisons :
\begin{itemize}
  \item \textbf{Échelle commune :} Sans elle, une variable à grande amplitude comme \texttt{date\_normalized} (max $= 2~782~661$ secondes) ou \texttt{ip\_interval\_std} (max $= 1~885~947$) dominerait entièrement les calculs de distance ou de variance, reléguant au second plan les variables à faible amplitude comme \texttt{ip\_request\_rate} (max $= 79$) ou \texttt{method\_encoded} (max $\approx 1$).
  \item \textbf{Compatibilité avec les modèles :} Les algorithmes sensibles aux échelles --- PCA, Isolation Forest, k-means, ou encore les réseaux de neurones --- requièrent des données centrées-réduites pour converger correctement et éviter qu'une seule variable ne conditionne l'ensemble de l'analyse.
  \item \textbf{Préparation du jeu de données final :} Appliquée après la transformation logarithmique, la standardisation garantit que toutes les variables, quelle que soit leur distribution d'origine, soient comparables en vue de la modélisation. Au terme du pipeline, le jeu de données final compte \textbf{41~433 lignes et 62 colonnes} (12 colonnes brutes enrichies par l'ingénierie des caractéristiques, l'encodage et la transformation logarithmique), dont 57 colonnes numériques standardisées.
\end{itemize}

\end{document}
